\begin{proposition}[4.2.4]
\label{IV.2.4.2.4-fr}
Sean $k$ un cuerpo, $X$ e $Y$ dos $k$-preesquemas localmente noetherianos
tales que $X\times_k Y$ sea localmente noetheriano. Supongamos además que $X$ e
$Y$ son íntegros. Entonces~:
\begin{enumerate}
\item[\textup{(i)}] $X\times_k Y$ carece de ciclos primos asociados inmersos~; cada una de las
componentes irreducibles de $X\times_k Y$ domina a $X$ y a $Y$, y el conjunto de
estas componentes está en correspondencia biunívoca con el conjunto de las
componentes irreducibles de
$\operatorname{Spec}(R(X)\otimes_k R(Y))$ \textup{(es decir,} el conjunto de los
ideales primos mínimos de $R(X)\otimes_k R(Y)$\textup{)}, donde $R(X)$ y $R(Y)$ son,
respectivamente, los cuerpos de funciones racionales de $X$ y de $Y$.
\item[\textup{(ii)}] Si un punto maximal $z$ de $X\times_k Y$ se identifica con un ideal
primo mínimo $\mathfrak p$ de $R(X)\otimes_k R(Y)$, el anillo local
$\mathcal O_{X\times_k Y,z}$ es isomorfo al anillo de fracciones
$(R(X)\otimes_k R(Y))_{\mathfrak p}$. En particular, si $R(X)$ o $R(Y)$ es
separable sobre $k$, $X\times_k Y$ es reducido.

\oldpage[IV-2]{56}

\item[\textup{(iii)}] Si además $X$ e $Y$ son localmente de tipo finito sobre $k$,
toda componente irreducible de $X\times_k Y$ tiene dimensión
$\dim(X)+\dim(Y)$.
\end{enumerate}
\end{proposition}

\begin{proof}
La afirmación (iii) se deduce de (4.2.1.2), teniendo en cuenta (i) y (ii). Para
demostrar (ii), es posible evidentemente limitarse al caso en que
$X=\operatorname{Spec}(A)$, $Y=\operatorname{Spec}(B)$ son afines, siendo $A$ y
$B$, por tanto, anillos íntegros con cuerpos de fracciones respectivos
$K=R(X)$ y $L=R(Y)$~; la afirmación (i) muestra que todo ideal primo mínimo
$\mathfrak q$ de $A\otimes_k B$ es la traza sobre $A\otimes_k B$ de un ideal
primo mínimo $\mathfrak p$ de $K\otimes_k L$~; como $K\otimes_k L$ es un
anillo de fracciones de $A\otimes_k B$, el isomorfismo de los anillos
$(A\otimes_k B)_{\mathfrak q}$ y $(K\otimes_k L)_{\mathfrak p}$ se deduce de
$(\mathbf{0}_{\mathrm I},1.2.5)$. Finalmente, si, por ejemplo, $R(Y)$ es separable sobre $k$,
se sabe que el anillo $R(X)\otimes_k R(Y)$ es reducido (Bourbaki,
\emph{Alg.}, chap.~VIII, \S~7, n\textsuperscript{o}~3, th.~1)~; lo mismo ocurre
con los anillos locales de los puntos maximales de $X\times_k Y$. Se deduce
que $X\times_k Y$ es reducido (3.2.1).
\end{proof}

\begin{proposition}[4.2.5]
\label{IV.2.4.2.5-fr}
Sean $k$ un cuerpo, $X$ e $Y$ $k$-preesquemas localmente noetherianos,
$\mathcal F$ \textup{(resp. $\mathcal G$)} un $\mathcal O_X$-módulo \textup{(resp.} un
$\mathcal O_Y$-módulo\textup{)} cuasicoherente. Sea $(Z'_\lambda)$ \textup{(resp.
$(Z''_\mu)$)} la familia de los ciclos primos asociados a $\mathcal F$ \textup{(resp.
$\mathcal G$)}, y designemos también por $Z'_\lambda$ \textup{(resp. $Z''_\mu$)} el
subpreesquema reducido de $X$ \textup{(resp. $Y$)} cuyo espacio subyacente es
$Z'_\lambda$ \textup{(resp. $Z''_\mu$)}. Entonces, si $Z'_\lambda\times_k Z''_\mu$
es localmente noetheriano, las componentes irreducibles $Z_{\lambda\mu\nu}$
de $Z'_\lambda\times_k Z''_\mu$ dominan a $Z'_\lambda$ y a $Z''_\mu$, y
$(Z_{\lambda\mu\nu})$ es la familia de los distintos ciclos primos asociados a
$\mathcal F\otimes_k\mathcal G$.
\end{proposition}

\begin{proof}
Basta aplicar (4.2.4) al producto
$\operatorname{Spec}(\boldsymbol{k}(x))\times_k\operatorname{Spec}(\boldsymbol{k}(y))$.
\end{proof}

En particular~:

\begin{corollary}[4.2.6]
\label{IV.2.4.2.6-fr}
Sean $k$ un cuerpo, $X$, $Y$ dos $k$-preesquemas localmente noetherianos tales
que $X\times_k Y$ sea localmente noetheriano. Sea $(Z'_\lambda)$ \textup{(resp.
$(Z''_\mu)$)} la familia de los subpreesquemas reducidos de $X$ \textup{(resp. $Y$)} que
tienen como espacios subyacentes las componentes irreducibles de $X$ \textup{(resp.
$Y$)}. Entonces las componentes irreducibles $Z_{\lambda\mu\nu}$ de
$Z'_\lambda\times_k Z''_\mu$ dominan a $Z'_\lambda$ y a $Z''_\mu$, y
$(Z_{\lambda\mu\nu})$ es la familia de las componentes irreducibles de
$X\times_k Y$.
\end{corollary}

\begin{proof}
En efecto, es posible limitarse al caso en que $X$ e $Y$ son reducidos
(\textbf I, 5.1.8)~; las componentes irreducibles de $X$ (resp. $Y$) son entonces los
ciclos primos asociados a $\mathcal O_X$ (resp. $\mathcal O_Y$) (3.2.1). Apliquemos (4.2.5)
a $\mathcal F=\mathcal O_X$ y $\mathcal G=\mathcal O_Y$, observando que por
definición $\mathcal O_X\otimes_k\mathcal O_Y=\mathcal O_{X\times_k Y}$~; el
corolario se deduce de que $\mathcal O_X\otimes_k\mathcal O_Y$ no tiene ciclos
primos asociados inmersos, pues lo mismo sucede con $\mathcal O_X$ y
$\mathcal O_Y$ por hipótesis (4.2.3).
\end{proof}

Apliquemos los resultados precedentes al caso en que $Y$ es el espectro de una
\emph{extensión} $K$ de $k$~:
